\documentclass[12pt]{article}

% \linespread{1.35} % for line distances, not change this

\usepackage[colorlinks=true, urlcolor=blue, linkcolor=blue]{hyperref}
\usepackage{amsmath, amssymb, amsfonts, amsthm}
\usepackage{breqn} % math related
\usepackage{paralist} % to put items in line
\usepackage[most]{tcolorbox}
\usepackage{fontawesome5}
\usepackage{xcolor}
\tcbuselibrary{skins,breakable,theorems}
\usepackage{xepersian}

\definecolor{QuestionBlue}{RGB}{30,90,170}
\definecolor{QuestionBack}{RGB}{245,248,255}
\definecolor{MyGreen}{RGB}{40,120,60}
\definecolor{MyGreenBack}{RGB}{243,250,243}

\emergencystretch=1em

\newtcbtheorem[
    auto counter,
    number within=section
]{question}{سؤال}{
    enhanced,
    breakable,
    colback=QuestionBack,
    colframe=QuestionBlue,
    coltitle=white,
    fonttitle=\bfseries,
    title filled,
    attach boxed title to top right={
        xshift=-2mm,
        yshift=-2mm
    },
    boxed title style={
        colback=QuestionBlue,
        colframe=QuestionBlue,
        arc=2mm
    },
    arc=2mm,
    boxrule=0.8pt,
    left=2mm,
    right=2mm,
    top=2mm,
    bottom=2mm,
    before skip=12pt,
    after skip=12pt,
}{q}

\newtcolorbox{greenbox}[1]{
    enhanced,
    breakable,
    colback=MyGreenBack,
    colframe=MyGreen,
    coltitle=white,
    fonttitle=\bfseries,
    title={#1},
    title filled,
    attach boxed title to top right={
        xshift=-4mm,
        yshift=-2mm
    },
    boxed title style={
        colback=MyGreen,
        colframe=MyGreen,
        arc=2mm
    },
    arc=2mm,
    boxrule=0.8pt,
    left=2mm,
    right=2mm,
    top=2mm,
    bottom=2mm,
    before skip=10pt,
    after skip=10pt
}

\newenvironment{solution}
{
    \par\medskip
    \hrule
    \smallskip
    \noindent\textbf{راه‌حل.}
    \par\smallskip
}
{
    \par\smallskip
    \smallskip
    \hrule
    \medskip
}



% \settextfont{XB Roya}
\settextfont{Sharif 1.2 Regular}

\begin{document}

\begin{titlepage}

\centering

\vspace*{2cm}

{\Large\textbf{دانشگاه صنعتی شریف}}

\vspace{0.5cm}

{\large دانشکده مهندسی کامپیوتر}

\vspace{2cm}

{\Huge\bfseries گزارش پروژه اول}

\vspace{0.5cm}

{\Huge تاب‌آوری راهبردی تحت عدم‌قطعیت}

\vspace{2.5cm}

{\large
\begin{tabular}{rc}
    \textbf{درس:} & نظریه بازی‌ها\\[0.3cm]
    \textbf{استاد:} & دکتر سمیرا حسین‌قربان\\[0.3cm]
    \textbf{دانشجو:} & ایلیا یزدانی
\end{tabular}
}

\vfill

{\large تابستان ۱۴۰۵}

\end{titlepage}

\tableofcontents

\newpage

\section{مقدمه: فهم مدل}
در این بخش، اجزای اصلی این مدل را بررسی می‌کنیم:

\begin{itemize}
    \item {
        \textbf{بازیگران اصلی:}
        دو بازیکن در این بازی حضور دارند. ایران که با 
        \lr{I}
        نمایش داده می‌شود و آمریکا که با
        \lr{U}
        نمایش می‌دهیم.
    }
    \item {
        \textbf{نوع خصوصی هر بازیکن:}
        هر بازیکن \lr{i} یک نوع ساختاری خصوصی به صورت
        \lr{$\theta_i = (\bar{\rho}_i, \bar{m}_i)$}
        دارد که در آغاز بازی تعیین می‌شود و در طول بازی ثابت می‌ماند؛
        \lr{$\bar{\rho}_i$}
        عزم پایه و
        \lr{$\bar{m}_i$}
        ظرفیت پایه مدیریت هزینه است. باور پیشین درباره نوع هر بازیکن نیز با
        \lr{$p_i(\theta_i)$}
        نمایش داده می‌شود.
    }
    \item {
        \textbf{مفهوم تاب‌آوری و تفاوت آن برای دو طرف:}
        تاب‌آوری با
        $e_i^t$
        نمایش داده می‌شود؛ متغیری پنهان و پویا که سطح تاب‌آوری بازیکن در مرحله
        \lr{$t$}
        است و از نوع ساختاری سرچشمه می‌گیرد
        (\lr{$e_i^0 = e_i^0(\theta_i)$}).
        برای ایران، تاب‌آوری یعنی جذب فشار خارجی بدون آن‌که این فشار به بی‌ثباتی داخلی تبدیل شود
        (مدیریت تحریم‌ها، حفظ نظم داخلی و انسجام نخبگان)؛
        و برای آمریکا، یعنی تداوم فشار بر ایران بدون خستگی داخلی، انزوای دیپلماتیک،
        گسترش بیش‌ازحد نظامی یا تشدیدِ مهارنشدنی تنش.
    }
    \item {
        \textbf{اقدامات و سیگنال‌ها:}
        اقدام کامل بازیکن \lr{i} در مرحله \lr{$t$} برابر است با
        \lr{$a_i^t = (\sigma_i^t, d_i^t)$}.
        مؤلفه
        \lr{$\sigma_i^t$}
        بخش سیگنالیِ مشاهده‌پذیر اقدام است
        (وضعیت نظامی، پیام دیپلماتیک، خویشتنداری، تشدید تنش، تعهد عمومی یا عملکرد زیر فشار).
        تاریخ عمومی اقدامات مشاهده‌شده با
        \lr{$h^t$}
        و مجموعه بازیگران فعال با
        \lr{$M(h^t)$}
        نمایش داده می‌شوند؛ به‌گونه‌ای که بسته به تاریخ، گاهی یک بازیکن زودتر حرکت می‌کند و گاهی هر دو هم‌زمان.
    }
    \item {
        \textbf{تصمیم ادامه یا خروج:}
        مؤلفه
        \lr{$d_i^t \in \{C, E\}$}
        تصمیم ادامه است:
        \begin{itemize}
            \item 
            ادامه تقابل 
            (\lr{C})

            \item 
             خروج/کاهش تنش/مصالحه/تغییر راهبرد (\lr{E}).
        \end{itemize}
        زمان توقف هر بازیکن نیز به‌صورت
        \lr{$\tau_i = \inf\{t : d_i^t = E\}$}
        تعریف می‌شود.
    }
    \item {
        \textbf{هزینه مؤثر ادامه:}
        هزینه جاری مؤثر ادامه تقابل برابر است با:
\[
g_i(t,h^t;\theta_i,e_i^t) = \dfrac{c_i(t,h^t) + a_i^{aud}(t,h^t) + r_i(t,h^t)}{m_i^t}
\]
        که در آن
        $c_i$
        هزینه مادی-اقتصادی،
        $a_i^{aud}$
        هزینه مخاطب داخلی و
        $r_i$
        هزینه ریسک تشدید است؛ مخرج کسر،
        $m_i^t$،
        ظرفیت جاری مدیریت هزینه است که توان جذب، توزیع و مدیریت سیاسی هزینه‌های خام را می‌سنجد.
        هزینه مؤثر انباشته تا مرحله
        \lr{$t$}
        نیز با
        \lr{$G_i(t;\cdot)$}
        نمایش داده می‌شود. نکته کلیدی این است که سطح یکسانی از فشار خام، بسته به ظرفیت مدیریت هزینه، هزینه مؤثر متفاوتی تولید می‌کند.
    }
    \item {
        \textbf{هزینه سیگنال‌دهی:}
        سیگنال‌ها پرهزینه‌اند و هزینه آن‌ها برابر است با:
        \[
            \varphi_i(\sigma_i^s;\theta_i,e_i^s,h^s) = \varphi_i^{prod} + \varphi_i^{commit}،
        \]
        یعنی هزینه تولید و نگهد سیگنال به‌علاوه هزینه تعهد و مخاطب.
        سیگنال قوی‌تر پرهزینه‌تر است، اما برای بازیکنی با تاب‌آوری،
         عزم یا ظرفیت مدیریت بیشتر، کم‌هزینه‌تر است.
           }
    \item {
        \textbf{ارزش ادامه دادن تا عقب‌نشینی طرف مقابل:}
        اگر بازیکن تا خروج حریف در بازی بماند، ارزش
        \lr{$B_i(e_i^t,\theta_i)$}
        را کسب می‌کند که در تاب‌آوری جاری و عزم پایه صعودی است.
    }
    \item {
        \textbf{ارزش خروج، مصالحه یا کاهش تنش:}
        ارزش خروج با
        \lr{$X_i(t,h^t)$}
        نمایش داده می‌شود و پویاست، یعنی به تاریخ عمومی بستگی دارد.
        خروج لزوماً به معنای شکست نیست؛ می‌تواند مصالحه، آتش‌بس یا تغییر راهبرد باشد.
    }
    \item {
        \textbf{باورها و به‌روزرسانی آن‌ها:}
        باور پسین بازیکن \lr{j} درباره نوع و تاب‌آوری جاری حریف با
        \lr{$\mu_j^t(\theta_i, e_i^t \mid h^t)$}
        نمایش داده می‌شود. چون تاب‌آوری مستقیم مشاهده نمی‌شود، هر طرف آن را از اقدامات،
        سیگنال‌ها و عملکرد طرف مقابل زیر فشار استنتاج می‌کند و در صورت امکان با قاعده بیز به‌روزرسانی می‌کند؛
        از این رو، این تقابل نه‌فقط یک مسابقه فشار مادی، بلکه یک مسابقه تفسیر است.
    }
\end{itemize}

در نهایت، مطلوبیت تحقق‌یافته هر بازیکن (\lr{$u_i$}) ترکیب همین اجزاست:
اگر بازیکن حریف را وادار به عقب‌نشینی کند (\lr{$\tau_i > \tau_j$}) ارزش
$B_i$
و اگر زودتر (یا هم‌زمان) خارج شود ارزش
$X_i$
را دریافت می‌کند؛ اما در هر دو حالت، هزینه‌های مؤثر انباشته
$G_i$
و هزینه‌های سیگنال‌دهی را تا لحظه توقف می‌پردازد.
بنابراین تصمیم مرکزی در هر مرحله این است که آیا ارزش انتظاری ادامه از ارزش خروج بیشتر است یا نه،
و تقابل تا زمانی تداوم می‌یابد که هر دو طرف ادامه را بر خروج ترجیح دهند.

\section{تحلیل نظری}

\subsection{تحلیل تصمیم ادامه یا خروج}

در این بخش نشان می‌دهیم که چگونه منطق بازی تقابل فرسایشی \footnote{War of Attrition}
از مدل داده‌شده استخراج می‌شود.
در هر مرحله $t$،
 هر بازیکن فعال 
 $i \in M(h^t)$
  پس از مشاهده تاریخ عمومی 
  $h^t$،
با دو گزینه کلی روبه‌روست:

\begin{itemize}
    \item {
        \textbf{ادامه تقابل:} 
        $d_i^t = C$
        ؛ یعنی ماندن در تقابل و
         پذیرفتن 
        هزینه‌های مؤثر و سیگنال‌دهی در آن مرحله؛
    }
    \item {
        \textbf{خروج:} 
        $d_i^t = E$
        ؛ یعنی بسته به تاریخ، خروج، کاهش تنش، مصالحه، آتش‌بس یا تغییر راهبرد،
        که ارزش خروج  $X_i(t,h^t)$ را تحقق می‌بخشد.
    }
\end{itemize}

چون بازی چندمرحله‌ای با اقدامات مشاهده‌شده و اطلاعات ناقص است، این تصمیم باید
عقلانی
 باشد: پس از هر تاریخ عمومی، انتخاب بازیکن باید با توجه به باورش درباره
نوع و تاب‌آوری حریف، یعنی 
$\mu_i^t(\theta_j, e_j^t \mid h^t)$،
 بهینه باشد.

\subsubsection*{شرط عقلانیت برای ادامه}

فرض کنید بازیکن $i$ در مرحله $t$ 
ماندن در بازی را در نظر می‌گیرد. پیامد نهایی او به مقایسه
زمان توقف خودش 
$\tau_i = \inf\{t : d_i^t = E\}$
 با زمان توقف حریف 
 $\tau_j$ 
 بستگی دارد:
اگر $\tau_i > \tau_j$،
 ارزش عقب‌نشینی حریف 
 $B_i(e_i^{\tau_j}, \theta_i)$
  را می‌گیرد
و اگر 
$\tau_i \le \tau_j$، 
ارزش خروج 
$X_i(\tau_i, h^{\tau_i})$
 را؛ و در هر دو حالت،
هزینه‌های مؤثر و سیگنال‌دهی انباشته تا پایان بازی از آن کسر می‌شود.
بنابراین ارزش مورد انتظار ادامه برابر است با:

\begin{align*}
    V_i^C(t,h^t) = 
    \mathbb{E}\Big[ \\
        & B_i(e_i^{\tau_j},\theta_i)\,\mathbf{1}_{\tau_i > \tau_j} \\
        & + X_i(\tau_i,h^{\tau_i})\,\mathbf{1}_{\tau_i \le \tau_j} \\
        & - \sum_{s=t+1}^{\tau} g_i(s,h^s;\theta_i,e_i^s)  \\
        & - \sum_{s=t+1}^{\tau} \varphi_i(\sigma_i^s;\theta_i,e_i^s,h^s) \\
    & \;\Big|\; h^t,\theta_i,e_i^t,\mu_i^t \Big],
\end{align*}

که در آن $\tau = \min\{\tau_i,\tau_j\}$ افق هزینه‌پردازی است و انتظار بر 
اساس باور بازیکن
درباره نوع، تاب‌آوری و بازی آینده حریف 
(و در نتیجه درباره $\tau_j$)
 گرفته می‌شود.
 توجه کنید که 
 $\mathbf{1}_p$
 اگر اتفاق 
 $p$
 رخ دهد برابر یک، وگرنه $0$
 می شود.

از سوی دیگر، اگر بازیکن در همین مرحله خارج شود،

\begin{equation*}
    V_i^E(t,h^t) = X_i(t,h^t).
\end{equation*}

توجه کنید که هزینه مؤثر انباشته 
$G_i(t;\cdot)$
 و هزینه‌های سیگنال‌دهی
$\sum_{s=0}^{t}\varphi_i(\sigma_i^s)$
 تا مرحله $t$
در هر دو شاخه به‌طور یکسان ظاهر می‌شوند و لذا از مقایسه حذف شده‌اند؛
گذشته، تنها از طریق تاریخ عمومی $h^t$،
 ارزش خروج پویا و باورهای به‌روزشده بر تصمیم جاری اثر می‌گذارد.
بنابراین، ادامه زمانی انتخاب می‌شود که منفعت مورد انتظار از ماندن در بازی، پس از کسر
هزینه‌های مؤثر و هزینه‌های سیگنال‌دهی، از ارزش خروج بیشتر باشد:

\begin{equation}
    \label{eq:continue}
    d_i^t = C \quad\Longleftrightarrow\quad
    V_i^C(t,h^t) \;\ge\; X_i(t,h^t).
\end{equation}

این همان منطق جنگ فرسایشی است که در این مدل نهفته است؛
 هر بازیکن تا زمانی در مسابقه پرهزینه
می‌ماند که ارزش انتظاری پیشی‌گرفتن بر حریف از ارزش ترک بازی 
بیشتر باشد؛ با این تفاوت که در اینجا
هزینه‌ها نه خام بلکه 
مؤثر
 اند (بر ظرفیت مدیریت هزینه $m_i^t$ تقسیم می‌شوند)،
ارزش خروج نه عددی ثابت بلکه 
  وابسته به تاریخ
 است.

از نامساوی~\eqref{eq:continue} روشن می‌شود که چه عواملی ادامه را محتمل‌تر می‌کنند:

\begin{itemize}
    \item احتمال ذهنی بالاترِ خروج زودهنگام حریف 
    (یعنی $\mu_i^t$ احتمال بیشتری به $e_j^t$ یا $\bar{\rho}_j$ پایین بدهد)، 
    که عبارت $B_i$ را بالا می‌برد؛
    \item ظرفیت مدیریت هزینه بالاتر 
    $m_i^t$، 
    که هزینه جاری مؤثر $g_i$
     را کوچک و انتظار را ارزان می‌کند؛
    \item تاب‌آوری جاری $e_i^t$
     و عزم پایه 
     $\bar{\rho}_i$
      بالاتر، که 
      $B_i$ 
      را افزایش می‌دهند؛
    \item  سرانجام، ارزش خروج جاری $X_i(t,h^t)$ پایین‌تر.
\end{itemize}

\subsubsection*{چرا یک بازیکن ممکن است با وجود هزینه‌بر بودن، ادامه دهد؟}

بر اساس همان نامساوی، مدل نشان می‌دهد که ادامه پرهزینه لزوماً رفتار غیرعقلانی نیست، و
گاهی هزینه خروج از ادامه‌دادن پرهزینه‌تر 
می‌شود. دلایل اصلی این موضوع در این مدل عبارتند از:

\begin{enumerate}
    \item \textbf{در دسترس‌نبودن سیاسیِ خروج.}
    اگر هزینه مخاطب داخلی، مصالحه یا کاهش تنش را 
    تحقیرآمیز
     یا ناممکن کرده باشد و هیچ مسیری با
         حفظ‌آبرو، میانجیگری معتبر یا گام متقابلی وجود نداشته باشد،
     $X_i(t,h^t)$
      بسیار پایین می‌آید؛
    در این حالت حتی با هزینه مؤثر بالا، نامساوی~\eqref{eq:continue} 
    برقرار می‌ماند:
    بازیکن نه به این دلیل ادامه می‌دهد که پیروزی محتمل است، بلکه
     به این دلیل که خروج از نظر سیاسی پرهزینه‌تر است.
    \item \textbf{انتظار خروج پیش‌تر حریف.}
    حتی اگر هزینه جاری $g_i$ 
    بالا باشد، چنانچه بازیکن باور داشته‌باشد دشمنش به انتهای
    تاب‌آوری خودش نزدیک شده و به زودی از بازی خارج می‌شود، آنگاه طبق منطق بازی
    فرسایشی ماندن در بازی در نهایت سودمندتر خواهد بود.
    \item \textbf{خودِ ادامه، یک سیگنال است.}
    چون تاب‌آوری مستقیم مشاهده نمی‌شود، حریف آن را از رفتار استنتاج می‌کند؛
     خروج پس از سیگنال‌های پرهزینه،
    تاب‌آوری پایین را آشکار می‌کند و اعتبار سرمایه‌گذاری‌شده را از بین می‌برد،
     اما ادامه، باور حریف به تاب‌آوری بالا
    را حفظ می‌کند و می‌تواند
     او را به خروج زودتر وادارد؛ یعنی انتخاب 
     $d_i^t = C$ بر باور آینده حریف
    $\mu_j^{t+1}$
     و در نتیجه توزیع 
     $\tau_j$
      اثر می‌گذارد.
      
    \item \textbf{شکاف هزینه خام و هزینه مؤثر.}
    فشار خام 
    (تحریم، فشار نظامی و مانند آن) ممکن است زیاد باشد، 
    اما اگر 
    $m_i^t$ بالا باشد،
    هزینه مؤثر $g_i$ 
    قابل‌مدیریت می‌ماند؛ بنابراین ادامه‌ای که از بیرون پرهزینه دیده می‌شود،
    برای خود بازیکن عقلانی است.
\end{enumerate}

در نهایت توجه کنید که تقابل تا زمانی تداوم می‌یابد که هر دو
 بازیکن ادامه را بر خروج ترجیح دهند؛
از این رو، بن‌بست طولانی و پرهزینه نه نشانه رفتار غیرعقلانی، بلکه می‌تواند بر مسیر تعادلی
یک بازی فرسایشی با اطلاعات ناقص باشد.

\subsection{تحلیل هزینه مؤثر}

هزینه ادامه تقابل در این مدل تنها هزینه خام نیست،
 آنچه بر تصمیم ادامه یا خروج اثر می‌گذارد هزینه مؤثر
 است و هزینه مؤثر به ظرفیت مدیریت هزینه نیز بستگی دارد:

\begin{equation}
    \label{eq:geff}
    g_i(t,h^t;\theta_i,e_i^t) = \frac{c_i(t,h^t) + a_i^{aud}(t,h^t) + r_i(t,h^t)}{m_i^t},
    \quad m_i^t = m_i(\theta_i,e_i^t,h^t) > 0.
\end{equation}

صورت کسر سه جزء اصلی دارد: هزینه مادی و اقتصادی 
$c_i$،
 هزینه سیاسی و اجتماعی داخلی 
$a_i^{aud}$،
 و هزینه ریسک تشدید و از دست رفتن کنترل بحران 
 $r_i$. اما این هزینه‌ها پیش از ورود به تصمیم‌گیری بر ظرفیت مدیریت هزینه 
 $m_i^t$
  تقسیم می‌شوند؛ بنابراین دو طرف که با فشار خام
   مشابهی روبه‌رویند، لزوماً هزینه مؤثر مشابهی تجربه نمی‌کنند.
   حال چهار پرسش تحلیلی مطرح شده را بررسی می‌کنیم:

\subsubsection*{اثر افزایش هزینه خام بر تصمیم ادامه:}
        افزایش هزینه خام 
        (تحریم شدیدتر، استقرار نظامی بیشتر و \dots)
         جریان هزینه مؤثر $g_i$
          را بالا می‌برد، ارزش ادامه $V_i^C$
           را کاهش می‌دهد و نامساوی~\eqref{eq:continue}
            را به سمت خروج متمایل می‌کند. اما چون 
            $\partial g_i/\partial c_i = 1/m_i^t$
            ، شدت این اثر به ظرفیت مدیریت هزینه بستگی دارد:
             همان افزایش فشار خام برای بازیکنی با 
             $m_i^t$
         بالاتر، اثر بسیار کوچک‌تری بر تصمیم ادامه دارد.
        
\subsubsection*{اثر افزایش ظرفیت مدیریت هزینه:}
        چون $\partial g_i/\partial m_i^t < 0$،
         افزایش $m_i^t$
          هر سه جزء هزینه خام را به
           بار مؤثر کمتری تبدیل می‌کند؛ انتظار کم‌هزینه‌تر می‌شود و
           ادامه دادن محتمل‌تر. افزون بر این، طبق بخش سیگنال‌دهی مدل،
            $\partial \varphi_i/\partial \bar{m}_i < 0$
            ؛ یعنی ظرفیت بالاتر، سیگنال عزم را نیز
     ارزان‌تر و معتبرتر می‌کند.
     در نتیجه، افزایش مدیریت هزینه هم بر ادامه
     دادن و هم سیگنال‌دهی 
     اثر مثبت دارد.
        
\subsubsection*{چرا فشار خارجی لزوماً به کاهش تاب‌آوری نمی‌انجامد؟}
        زیرا تاب‌آوری معادل تحمل درد خام نیست؛ فشار تنها وقتی تاب‌آوری 
        $e_i^t$
    را فرسایش می‌دهد که
    به هزینه‌های داخلی غیرقابل‌مدیریت تبدیل شود. اگر 
    $m_i^t$ 
    بالا بماند، مخرج بزرگ کسر رابطه~\eqref{eq:geff} 
    افزایش صورت را خنثی می‌کند و هزینه مؤثر قابل‌مدیریت می‌ماند؛
    حتی عملکرد موفق زیر فشار، باور حریف 
    $\mu_j^t$
    درباره تاب‌آوری را به نفع بازیکن نگه می‌دارد. 
    پس فشار خارجی بدون شکاف در مدیریت داخلی،
        لزوماً
        باعث کاهش تاب‌آوری و 
        نزدیک شدن بازیکن به خروج
        نمی‌شود.
        
\subsubsection*{چرا ضعف در مدیریت داخلی هزینه، اعتبار تاب‌آوری را کاهش می‌دهد؟}
    ضعف مدیریتی 
    (بی‌ثباتی اقتصادی، نارضایتی عمومی، شکاف نخبگان)
    هم $a_i^{aud}$ و هم $g_i$ 
    را بالا می‌برد و هم برای حریف مشاهده‌پذیر
    است؛ حریف از این عملکرد، باور 
    $\mu_j^t(\theta_i,e_i^t \mid h^t)$
    را به سمت نوع‌های ضعیف به‌روزرسانی می‌کند
    و به خروج حریف در زمان نزدیک باور پیدا می‌کند.
     در نتیجه 
    سیگنال‌های عزم او دیگر باورپذیر نخواهند بود و حریف
    ترغیب می‌شود فشار را ادامه دهد.
     بنابراین ضعف
 مدیریتی داخلی نه تنها 
 تاب‌آوری راهبردی را کاهش می‌دهد، بلکه ادامه دادن را نیز پرهزینه می‌کند.

\vspace{0.7cm}

بنابراین، در این مدل
 فشار به‌خودی‌خود مزیت راهبردی تولید نمی‌کند؛ اثر آن همیشه از
 فیلتر ظرفیت مدیریت هزینه عبور می‌کند و دو طرف با فشار خام
 مشابه می‌توانند هزینه‌های مؤثر بسیار متفاوتی تجربه کنند.

\subsection{تحلیل باورها و اطلاعات ناقص}

\subsection{تحلیل سیگنال‌دهی پرهزینه}

\subsection{تحلیل تعارض اعتبار و انعطاف پذیری}

\subsection{تحلیل ارزش خروج}

\subsection{تحلیل تعادلی}

\subsection{گزاره تحلیلی}

\section{پیاده‌سازی (انشاالله خالی بماند)}

\section{تفسیر نتایج}

\section{محدودیت‌های مدل}


\end{document}